\chapter{Trabalhos relacionados}
\label{chap:trabalhos-relacionados}

Esta dissertação fica na interseção de três linhas de pesquisa: a que relaciona a poluição por plástico à população humana (\autoref{sec:rel-plastico}), a que define cidades a partir de dados em grade e estuda a sua distribuição de tamanhos (\autoref{sec:rel-cidades}) e a que aplica leis de escala a fenômenos urbanos (\autoref{sec:rel-escala}). A \autoref{sec:rel-posicionamento} resume as diferenças entre este trabalho e os anteriores.

\section{Plástico, rios e população}
\label{sec:rel-plastico}

\textbf{Evidência de campo.} \citeonline{browne2011} amostraram sedimentos de 18 praias em seis continentes e encontraram mais microplástico nas praias de regiões densamente povoadas; a maior parte das partículas eram fibras, compatíveis com a descarga de efluentes de máquinas de lavar. \citeonline{yonkos2014} mediram microplástico na superfície de quatro rios estuarinos da Baía de Chesapeake, nos Estados Unidos, e encontraram concentrações positivamente correlacionadas com a densidade populacional e com a proporção de área urbana e suburbana das bacias. Esses dois estudos são as evidências de campo mais diretas de que mais gente significa mais microplástico, mas têm alcance limitado: o primeiro compara praias sem delimitar a área que drena para cada uma, e o segundo compara apenas quatro bacias vizinhas.

\textbf{\textit{Pellets} plásticos.} \citeonline{tunnell2020} apresentaram os primeiros resultados do \textit{Nurdle Patrol}, com coletas de voluntários em praias de todo o Golfo do México. As maiores contagens ficaram na costa do Texas, região que concentra a indústria petroquímica e de produção de resinas, o que indica que, para os \textit{pellets}, a proximidade de fontes industriais pode pesar tanto quanto a população. \citeonline{karlsson2018} chegaram a conclusão semelhante ao medir a perda de \textit{pellets} ao redor de uma fábrica de polietileno na Suécia e estimar que milhões de grânulos por ano escapavam para o ambiente a partir de um único sítio industrial.

\textbf{Modelos de emissão por rios.} Os modelos globais de \citeonline{jambeck2015}, \citeonline{lebreton2017} e \citeonline{meijer2021}, descritos na \autoref{sec:microplastico}, partem da população e da gestão de resíduos para estimar a carga de plástico de cada país ou bacia. \citeonline{schmidt2017} calibraram um modelo desse tipo com medições de carga fluvial e concluíram que os dez rios com maior carga respondem pela maior parte do plástico exportado ao mar. Nesses trabalhos, a população entra como variável explicativa, e as medições servem para calibrar os coeficientes. Esta dissertação faz o caminho inverso: parte de medições independentes e pergunta se elas se organizam segundo a população das bacias.

\section{Cidades definidas a partir de dados em grade}
\label{sec:rel-cidades}

\citeonline{rozenfeld2008} introduziram o CCA e o aplicaram a dados de população em grade dos Estados Unidos, do Reino Unido e da África, mostrando que a taxa média de crescimento populacional das manchas urbanas não depende do tamanho, como prevê a lei de Gibrat, mas que a dispersão dessas taxas decai com o tamanho segundo uma lei de potência. \citeonline{rozenfeld2011} usaram o CCA com dados de alta resolução dos Estados Unidos e do Reino Unido para mostrar que tanto a população quanto a área das manchas seguem a lei de Zipf acima de um tamanho mínimo, em contraste com o resultado de \citeonline{eeckhout2004}, que encontrou uma distribuição lognormal usando lugares censitários. A diferença entre os dois resultados vem justamente da definição de cidade: os lugares censitários incluem muitas unidades pequenas que o CCA agruparia ou descartaria. \citeonline{jiang2011} chegaram à lei de Zipf nos Estados Unidos por outro caminho, definindo ``cidades naturais'' a partir da densidade de cruzamentos da malha viária.

Nesses trabalhos, os parâmetros do CCA são fixados a partir de valores típicos ou de critérios externos, como o limiar de área urbana do censo, e a lei de Zipf é o resultado a ser verificado. Aqui a lógica é inversa: a lei de Zipf é usada como critério para escolher os parâmetros, e o resultado de interesse é outro, a relação entre população e plástico. Essa inversão só é legítima porque a lei de Zipf é uma propriedade das cidades, independente do plástico, e porque o critério é definido antes de olhar para a variável de resposta.

\section{Leis de escala em fenômenos urbanos}
\label{sec:rel-escala}

\citeonline{bettencourt2007} estabeleceram o arcabouço de leis de escala urbana descrito na \autoref{sec:escala}. \citeonline{oliveira2014} aplicaram o CCA aos Estados Unidos e encontraram uma relação superlinear entre a população das manchas urbanas e as suas emissões de CO$_2$, com $\beta \approx 1{,}46$, concluindo que cidades grandes são menos eficientes em emissões por habitante. \citeonline{melo2014} encontraram relações de escala entre população e suicídios nos municípios brasileiros, e \citeonline{caminha2017} mostraram que, dentro de uma cidade, os crimes contra o patrimônio escalam de forma superlinear com a população flutuante e não com a residente. Este último resultado é relevante aqui por um motivo metodológico: a escolha de qual população medir (residente, flutuante, da cidade, da bacia) muda a relação encontrada.

\citeonline{arcaute2015} e \citeonline{cottineau2017} mostraram que os expoentes de escala são sensíveis à definição das unidades urbanas e alertaram para o risco de escolher a definição que produz o resultado esperado. Nesta dissertação, a unidade de análise não é a cidade, mas a bacia hidrográfica, que é definida pelo relevo e não pela população. Mesmo assim, as manchas urbanas entram na análise, e por isso os seus parâmetros são fixados por um critério independente.

\section{Posicionamento deste trabalho}
\label{sec:rel-posicionamento}

O Quadro~\ref{quadro:relacionados} compara esta dissertação com os trabalhos mais próximos. Nenhum dos trabalhos encontrados combina, numa mesma análise, (i) bacias delineadas a partir do relevo com todas as saídas costeiras, (ii) cidades definidas sem fronteiras administrativas e com parâmetros calibrados de forma independente, e (iii) medições de campo de microplástico associadas à foz de cada bacia, numa escala continental.

\begin{quadro}[ht!]
	\centering
	\caption{\label{quadro:relacionados}Comparação entre este trabalho e os trabalhos relacionados mais próximos}
	\small
	\begin{tabular}{|p{2.6cm}|p{2.6cm}|p{3.2cm}|p{3.4cm}|}
		\hline
		\textbf{Trabalho} & \textbf{Unidade de análise} & \textbf{Medida de população} & \textbf{Medida de plástico} \\
		\hline
		\citeonline{browne2011} & Praias (18) & Densidade populacional do entorno & Microplástico no sedimento (campo) \\
		\hline
		\citeonline{yonkos2014} & Bacias estuarinas (4) & Densidade populacional e uso do solo & Microplástico na superfície (campo) \\
		\hline
		\citeonline{lebreton2017} & Bacias globais & População e resíduos mal geridos & Carga modelada, calibrada com medições \\
		\hline
		\citeonline{meijer2021} & Bacias globais & População, resíduos, uso do solo & Carga modelada, calibrada com medições \\
		\hline
		\citeonline{tunnell2020} & Praias do Golfo do México & Não usa & \textit{Pellets} por 10 min (campo) \\
		\hline
		\citeonline{oliveira2014} & Manchas urbanas (CCA) & População das manchas & Não se aplica (CO$_2$) \\
		\hline
		\textbf{Esta dissertação} & Bacias costeiras (IPBA), 138 principais & População da bacia, com manchas urbanas (CCA) repartidas & \textit{Pellets} por 10 min na foz (campo) \\
		\hline
	\end{tabular}
	\par\medskip
	\footnotesize Fonte: elaborado pelo autor.
\end{quadro}
